\documentclass[12pt,a4paper]{report}
\usepackage[utf8]{inputenc}
\usepackage[T1]{fontenc}
\usepackage[french]{babel}
\usepackage{geometry}
\geometry{margin=2.5cm}
\usepackage{graphicx}
\usepackage{hyperref}
\usepackage{amsmath}
\usepackage{amssymb}
\usepackage{listings}
\usepackage[dvipsnames]{xcolor}
\usepackage{fancyhdr}
\usepackage{titlesec}
\usepackage{tocloft}
\usepackage{booktabs}
\usepackage{multirow}
\usepackage{float}
\usepackage{subcaption}

% Définition du thème vert et blanc professionnel
\definecolor{primarygreen}{RGB}{34, 139, 34}  % Vert forêt
\definecolor{secondarygreen}{RGB}{144, 238, 144}  % Vert clair
\definecolor{darkgreen}{RGB}{0, 100, 0}  % Vert foncé
\definecolor{lightgray}{RGB}{245, 245, 245}  % Gris très clair

% Configuration des hyperliens en vert
\hypersetup{
    colorlinks=true,
    linkcolor=primarygreen,
    citecolor=darkgreen,
    urlcolor=primarygreen,
    pdftitle={Rapport PFA - Matching CV/Offres d'emploi},
    pdfauthor={Ali},
    pdfsubject={Projet Final d'Année - Intelligence Artificielle},
    pdfkeywords={IA, Matching CV, Recrutement, NLP}
}

% Configuration pour les listings de code avec thème vert
\lstset{
    language=Python,
    basicstyle=\ttfamily\small,
    keywordstyle=\color{primarygreen}\bfseries,
    commentstyle=\color{darkgreen},
    stringstyle=\color{black},
    numbers=left,
    numberstyle=\tiny\color{primarygreen},
    frame=single,
    rulecolor=\color{primarygreen},
    breaklines=true,
    captionpos=b,
    backgroundcolor=\color{lightgray}
}

% En-têtes et pieds de page avec thème vert
\pagestyle{fancy}
\fancyhf{}
\fancyhead[L]{\color{primarygreen}\leftmark}
\fancyhead[R]{\color{primarygreen}\thepage}
\fancyfoot[C]{\color{darkgreen}\footnotesize{Rapport PFA - Matching CV/Offres d'emploi}}
\renewcommand{\headrulewidth}{1pt}
\renewcommand{\headrule}{\color{primarygreen}\hrule}
\renewcommand{\footrulewidth}{0.5pt}
\renewcommand{\footrule}{\color{secondarygreen}\hrule}

% Configuration de la table des matières
\setcounter{tocdepth}{2}
\setcounter{secnumdepth}{2}

% Style des titres de chapitre avec thème vert
\titleformat{\chapter}[display]
{\normalfont\color{primarygreen}\huge\bfseries}
{\color{primarygreen}\chaptertitlename\ \thechapter}{20pt}
{\color{primarygreen}\Huge}
\titlespacing*{\chapter}{0pt}{0pt}{40pt}

% Style des titres de section avec thème vert
\titleformat{\section}
{\normalfont\color{primarygreen}\Large\bfseries}
{\color{primarygreen}\thesection}{1em}{}
\titlespacing*{\section}{0pt}{3.5ex plus 1ex minus .2ex}{2.3ex plus .2ex}

% Style des titres de subsection
\titleformat{\subsection}
{\normalfont\color{darkgreen}\large\bfseries}
{\color{darkgreen}\thesubsection}{1em}{}
\titlespacing*{\subsection}{0pt}{3.25ex plus 1ex minus .2ex}{1.5ex plus .2ex}

% Configuration de la table des matières
\renewcommand{\cftchapfont}{\color{primarygreen}\bfseries}
\renewcommand{\cftsecfont}{\color{darkgreen}}
\renewcommand{\cftsubsecfont}{\color{black}}

% Page de titre personnalisée avec thème vert
\makeatletter
\def\@maketitle{
  \newpage
  \null
  \vskip 2em
  \begin{center}
    \color{primarygreen}
    {\Huge \bfseries \@title \par}
    \vskip 1.5em
    {\Large \bfseries \@author \par}
    \vskip 1em
    {\large \@date \par}
    \vskip 2em
    {\color{darkgreen}\rule{0.8\textwidth}{2pt}}
  \end{center}
  \par
  \vskip 1.5em
}
\makeatother

\title{Rapport de Projet Final d'Année (PFA) \\ \color{primarygreen}Moteur de Matching CV / Offres d'emploi \\ \color{darkgreen}Système de Matching Sémantique avec Intelligence Artificielle}
\author{\color{primarygreen}Ali}
\date{\color{darkgreen}Mai 2026}

\begin{document}

% Page de titre
\maketitle
\thispagestyle{empty}
\newpage

% Table des matières avec style vert
\tableofcontents
\newpage

% Liste des figures
\listoffigures
\newpage

% Liste des tableaux
\listoftables
\newpage

% Partie préliminaire
\part{Partie préliminaire}

% Dédicaces
\chapter{Dédicaces}
\label{chap:dedicaces}

{\color{primarygreen}\large\textbf{À mes chers parents,}} pour leur soutien inconditionnel, leurs sacrifices et leur encouragement constant tout au long de mon parcours académique et professionnel.

{\color{darkgreen}\large\textbf{À mes professeurs et encadrants,}} pour leur guidance experte et leur transmission des connaissances qui ont rendu ce projet possible.

{\color{primarygreen}\large\textbf{À mes collègues et amis,}} pour leur collaboration et leurs conseils précieux durant les phases de développement.

Que ce travail contribue à l'avancement de la recherche dans le domaine du recrutement assisté par l'IA.

% Remerciements
\chapter{Remerciements}
\label{chap:remerciements}

Je tiens à exprimer ma profonde gratitude à toutes les personnes qui ont contribué à la réalisation de ce projet :

\begin{itemize}
\item {\color{primarygreen}\textbf{Mon encadrant de PFA,}} pour ses conseils avisés et son suivi rigoureux tout au long du projet.
\item {\color{darkgreen}\textbf{L'équipe de développement,}} pour leur collaboration et leur expertise technique.
\item {\color{primarygreen}\textbf{Les contributeurs des datasets open-source utilisés,}} notamment les auteurs des datasets HuggingFace.
\item {\color{darkgreen}\textbf{La communauté open-source,}} particulièrement les développeurs de HuggingFace Transformers et Sentence Transformers.
\item {\color{primarygreen}\textbf{Mes professeurs de Master,}} pour les bases solides en intelligence artificielle et traitement du langage naturel.
\end{itemize}

Ce projet n'aurait pas été possible sans le soutien de ces personnes et institutions.

% Résumé
\chapter{Résumé}
\label{chap:resume}

Dans un contexte où les processus de recrutement à grande échelle deviennent de plus en plus complexes, ce rapport présente le développement d'un moteur de matching CV-Offres d'emploi utilisant les technologies de l'intelligence artificielle les plus avancées. Le système développé combine un modèle d'embeddings dense fine-tuné sur le domaine spécifique du recrutement, un graphe de proximité des compétences hand-curated, et un module d'extraction de signaux externes via les URLs publiques des candidats.

{\color{primarygreen}\large\textbf{L'approche proposée améliore significativement les performances par rapport aux méthodes traditionnelles,}} avec une corrélation Spearman de 0,78 contre 0,62 pour le modèle de base, représentant une amélioration relative de 25,8\%. Le système offre une explicabilité complète des résultats de matching, supporte le traitement multi-langues, et fournit une API REST production-ready avec un endpoint de matching batch adapté aux workflows des recruteurs.

Le rapport détaille l'ensemble du processus de développement, depuis l'analyse de marché et la conception système jusqu'à la validation expérimentale et le déploiement. Les résultats démontrent le potentiel des technologies d'IA pour révolutionner les processus de recrutement tout en maintenant une transparence et une équité essentielles.

{\color{darkgreen}\large\textbf{Mots-clés :}} Intelligence Artificielle, Traitement du Langage Naturel, Matching Sémantique, Recrutement, Embeddings, Fine-tuning, Explicabilité.

% Abstract
\chapter{Abstract}
\label{chap:abstract}

This report presents the development of a comprehensive CV-Job Matching Engine using state-of-the-art artificial intelligence technologies. The system combines a domain-adapted dense retrieval model (fine-tuned BAAI/bge-large-en-v1.5), a hand-curated skill proximity graph spanning 150+ technical skills across 11 technology clusters, and external URL signal extraction from public profiles (GitHub, LinkedIn, portfolios).

{\color{primarygreen}\large\textbf{The proposed approach significantly improves matching accuracy compared to traditional methods,}} achieving a Spearman rank correlation of 0.78 against human-labeled fit scores (up from 0.62 for the base model, representing a 25.8\% relative improvement). The system provides complete explainability of matching results, supports multi-language processing via automatic language detection and BGE-M3 routing, and delivers a production-ready REST API with batch matching capabilities suitable for recruiter workflows.

The report details the complete development process, from market analysis and system design through experimental validation and deployment. Results demonstrate the potential of AI technologies to revolutionize recruitment processes while maintaining essential transparency and fairness.

{\color{darkgreen}\large\textbf{Keywords:}} Artificial Intelligence, Natural Language Processing, Semantic Matching, Recruitment, Embeddings, Fine-tuning, Explainability.

% Introduction générale
\chapter{Introduction générale}
\label{chap:introduction}

\section{Contexte général}

Le recrutement constitue l'une des fonctions stratégiques les plus critiques pour les organisations modernes. Dans un environnement économique où la guerre des talents s'intensifie, les entreprises doivent traiter des volumes croissants de candidatures tout en maintenant des standards élevés de qualité et d'équité dans leurs processus de sélection.

{\color{primarygreen}\large\textbf{Selon les études récentes, un recruteur passe en moyenne 6 secondes sur l'examen initial d'un CV,}} un temps insuffisant pour une évaluation approfondie. Cette contrainte temporelle, combinée à la subjectivité inhérente aux jugements humains, conduit fréquemment à des biais de sélection et au passage à côté de candidats qualifiés dont les compétences ne sont pas explicitement mentionnées.

\section{Problème posé}

Le défi central adressé par ce projet est le développement d'un système automatisé capable de :
\begin{itemize}
\item {\color{darkgreen}Comprendre la sémantique profonde des descriptions de poste et des CV}
\item {\color{primarygreen}Évaluer objectivement l'adéquation candidat-poste}
\item {\color{darkgreen}Fournir des justifications transparentes des décisions}
\item {\color{primarygreen}Traiter efficacement des volumes importants de candidatures}
\end{itemize}

\section{Approche proposée}

Notre solution repose sur une architecture hybride combinant :
\begin{itemize}
\item {\color{primarygreen}Un modèle d'embeddings dense adapté au domaine du recrutement}
\item {\color{darkgreen}Un graphe de connaissances expertes sur les relations entre compétences}
\item {\color{primarygreen}L'extraction de signaux externes complémentaires}
\item {\color{darkgreen}Un système d'explicabilité basé sur des règles}
\end{itemize}

\section{Structure du rapport}

Ce rapport est structuré comme suit : le Chapitre~\ref{chap:contexte} présente le contexte général et la problématique ; le Chapitre~\ref{chap:marche} développe l'étude de marché et de faisabilité ; les Chapitres~\ref{chap:organisation} à~\ref{chap:realisation} détaillent respectivement l'organisation du projet, l'analyse et conception, la gestion des données, la modélisation IA et la réalisation technique ; enfin, le Chapitre~\ref{chap:resultats} présente les résultats et validations avant la conclusion générale.

% Chapitre 1
\part{Chapitre 1 : Contexte général et problématique}
\label{part:contexte}
\chapter{Contexte général et problématique}
\label{chap:contexte}

\section{Contexte du recrutement actuel}

\subsection{L'évolution du marché du travail}

Le marché du recrutement a connu une transformation profonde ces dernières années, caractérisée par :
\begin{itemize}
\item Une accélération de la rotation du personnel dans les secteurs technologiques
\item Une pénurie de talents qualifiés dans les domaines STEM (Science, Technology, Engineering, Mathematics)
\item L'émergence de nouveaux profils professionnels hybrides
\item Une digitalisation croissante des processus RH
\end{itemize}

\subsection{Les défis opérationnels des recruteurs}

Les recruteurs font face à plusieurs contraintes opérationnelles majeures :
\begin{itemize}
\item \textbf{Volume de candidatures :} Les postes attractifs génèrent souvent plusieurs centaines de CV
\item \textbf{Pression temporelle :} Délais de recrutement de plus en plus courts
\item \textbf{Complexité technique :} Évaluation de compétences techniques pointues
\item \textbf{Diversité des profils :} Candidats issus de parcours variés et non linéaires
\end{itemize}

\subsection{L'impact des biais cognitifs}

Les études en psychologie cognitive démontrent que les décisions humaines sont sujettes à de nombreux biais :
\begin{itemize}
\item Biais de confirmation : tendance à favoriser les candidats confirmant les préjugés
\item Biais d'ancrage : influence excessive de la première impression
\item Biais de similarité : préférence pour les profils similaires au sien
\item Biais de halo : généralisation d'une qualité à l'ensemble du profil
\end{itemize}

Ces biais sont particulièrement problématiques dans le contexte du recrutement où les décisions doivent être objectives et équitables.

\subsection{Le rôle croissant des technologies}

L'adoption des technologies dans le recrutement s'accélère :
\begin{itemize}
\item Systèmes de suivi des candidats (ATS - Applicant Tracking Systems)
\item Plateformes de recrutement en ligne
\item Outils d'analyse de CV automatisés
\item Intelligence artificielle pour le matching candidat-poste
\end{itemize}

\subsection{Évolution des attentes des candidats}

Les candidats modernes attendent :
\begin{itemize}
\item Des processus de recrutement transparents
\item Des retours rapides sur leurs candidatures
\item Des expériences utilisateur fluides
\item Une communication personnalisée
\end{itemize}

\section{Limites des méthodes traditionnelles}

\subsection{Les approches lexicales}

Les systèmes traditionnels de matching CV reposent principalement sur :
\begin{itemize}
\item La recherche de mots-clés exacts
\item Les systèmes de scoring basés sur des règles prédéfinies
\item Les bases de données relationnelles avec recherche booléenne
\end{itemize}

\subsubsection{Limitations fondamentales}

Ces approches souffrent de limitations fondamentales :
\begin{itemize}
\item \textbf{Absence de compréhension sémantique :} Incapacité à reconnaître les synonymes et termes équivalents
\item \textbf{Rigidité syntaxique :} Sensibilité à la formulation et à la structure des textes
\item \textbf{Contexte limité :} Manque de compréhension du contexte professionnel
\item \textbf{Évolutivité réduite :} Difficulté à adapter aux nouveaux domaines et terminologies
\end{itemize}

\subsubsection{Exemples de limitations}

Considérons l'exemple suivant :
\begin{itemize}
\item CV : "Développeur full-stack avec expertise en JavaScript et Node.js"
\item Offre : "Ingénieur logiciel maîtrisant JS et environnement serveur"
\end{itemize}

Un système lexical pourrait manquer ce matching évident en raison de l'absence de correspondance exacte des termes.

\subsection{Les solutions semi-automatisées}

Les ATS (Applicant Tracking Systems) modernes intègrent des fonctionnalités avancées mais restent limités par :
\begin{itemize}
\item Des algorithmes de ranking basés sur des scores composites arbitraires
\item Une dépendance forte aux données structurées saisies manuellement
\item Des capacités d'extraction d'informations limitées pour les CV non structurés
\end{itemize}

\subsubsection{Problèmes d'extraction d'informations}

L'extraction automatique d'informations depuis les CV présente plusieurs défis :
\begin{itemize}
\item Variabilité des formats et structures
\item Ambiguïtés linguistiques
\item Contexte culturel et sectoriel
\item Évolution des terminologies techniques
\end{itemize}

\subsection{L'échec des approches génériques}

L'application de modèles d'IA génériques au domaine du recrutement révèle plusieurs problèmes :
\begin{itemize}
\item Les embeddings pré-entraînés ne capturent pas les nuances du domaine professionnel
\item Absence de connaissances expertes sur les relations entre compétences
\item Manque d'explicabilité des décisions automatisées
\end{itemize}

\subsubsection{Le problème de la généralité}

Les modèles d'IA généralistes sont entraînés sur des corpus vastes mais génériques :
\begin{itemize}
\item Wikipedia, livres, articles de presse
\item Peu de contenu spécifique au recrutement
\item Terminologies techniques sous-représentées
\item Contextes professionnels mal couverts
\end{itemize}

\subsubsection{Exemple de limitation sémantique}

Prenons l'exemple des compétences transférables :
\begin{itemize}
\item Un développeur expérimenté en "Vue.js" peut facilement apprendre "React"
\item Un data scientist maîtrisant "scikit-learn" peut s'adapter à "TensorFlow"
\item Ces relations ne sont pas évidentes pour un modèle générique
\end{itemize}

\section{Problématique du matching CV / offres}

\subsection{Le défi sémantique}

Le problème central du matching CV-offres réside dans la nature fondamentalement sémantique de l'adéquation candidat-poste :
\begin{itemize}
\item Les compétences requises sont rarement listées explicitement
\item Les expériences passées doivent être interprétées dans leur contexte
\item Les soft skills et l'adaptabilité ne sont pas directement observables
\end{itemize}

\subsubsection{Complexité de l'évaluation des compétences}

L'évaluation des compétences techniques implique :
\begin{itemize}
\item Compréhension des niveaux de maîtrise (débutant, intermédiaire, expert)
\item Reconnaissance des compétences transférables
\item Évaluation de l'ancienneté et de la pertinence
\item Considération du contexte sectoriel
\end{itemize}

\subsubsection{Le rôle du contexte}

Le contexte joue un rôle crucial dans l'interprétation :
\begin{itemize}
\item "Python" dans un contexte data science vs développement web
\item "Leadership" dans une startup vs une grande entreprise
\item "International" pour un poste export vs local
\end{itemize}

\subsection{L'exigence d'explicabilité}

Dans un contexte réglementaire de plus en plus strict (RGPD, égalité des chances), les systèmes automatisés doivent :
\begin{itemize}
\item Justifier leurs décisions de manière compréhensible
\item Éviter les discriminations algorithmiques
\item Permettre la vérification et l'audit des processus
\end{itemize}

\subsubsection{Transparence algorithmique}

La transparence nécessite :
\begin{itemize}
\item Explication des critères de décision
\item Traçabilité des calculs effectués
\item Possibilité de contestation des résultats
\item Auditabilité des processus
\end{itemize}

\subsubsection{Éthique et équité}

Les considérations éthiques incluent :
\begin{itemize}
\item Évitement des biais discriminatoires
\item Respect de la diversité
\item Protection des données personnelles
\item Equité procédurale
\end{itemize}

\subsection{L'échelle et la performance}

Les solutions déployées doivent répondre aux contraintes opérationnelles réelles :
\begin{itemize}
\item Traitement de centaines de CV en quelques secondes
\item Intégration transparente dans les workflows existants
\item Robustesse face à la variabilité des formats et langues
\end{itemize}

\subsubsection{Contraintes de performance}

Les exigences de performance incluent :
\begin{itemize}
\item Latence acceptable (< 5 secondes pour un matching simple)
\item Débit élevé (100+ matchings/minute)
\item Disponibilité (> 99\%)
\item Évolutivité horizontale
\end{itemize}

\subsubsection{Défi de l'intégration}

L'intégration dans les workflows existants nécessite :
\begin{itemize}
\item APIs standardisées et documentées
\item Formats d'échange compatibles
\item Interfaces utilisateur intuitives
\item Formation minimale des utilisateurs
\end{itemize}

\section{Objectifs du projet}

\subsection{Objectifs généraux}

Ce projet vise à développer un système de matching CV-offres qui :
\begin{enumerate}
\item \textbf{Améliore la précision :} Surpasse les méthodes traditionnelles en termes de corrélation avec les jugements humains
\item \textbf{Garantit l'explicabilité :} Fournit des justifications claires et actionnables pour chaque décision
\item \textbf{Assure l'évolutivité :} Traite efficacement des volumes importants de données
\item \textbf{Maintient la flexibilité :} S'adapte à différents secteurs et types de poste
\end{enumerate}

\subsection{Objectifs spécifiques}

Les objectifs techniques détaillés incluent :
\begin{itemize}
\item Atteindre une corrélation Spearman ≥ 0,75 sur les scores de fit humain
\item Supporter au moins 5 langues majeures (Anglais, Français, Arabe, Espagnol, Allemand)
\item Traiter 100 CV contre 1 offre en moins de 30 secondes
\item Générer des explications couvrant ≥ 80\% des décisions de matching
\end{itemize}

\subsubsection{Métriques de performance cibles}

\begin{table}[H]
\centering
\caption{Objectifs de performance du système}
\label{tab:objectifs_performance}
\begin{tabular}{@{}lll@{}}
\toprule
Métrique & Valeur cible & Justification \\
\midrule
Corrélation Spearman & ≥ 0,75 & Surpasser les baselines existantes \\
Précision top-5 & ≥ 85\% & Qualité du ranking \\
Temps de réponse & < 30s & Utilisabilité opérationnelle \\
Couverture explicabilité & ≥ 80\% & Transparence requise \\
Disponibilité & > 99\% & Fiabilité production \\
\bottomrule
\end{tabular}
\end{table}

\subsection{Objectifs opérationnels}

Le système doit démontrer sa valeur dans des scénarios réels :
\begin{itemize}
\item Réduction du temps de screening de 80\%
\item Amélioration du taux de candidats qualifiés identifiés
\item Diminution des biais de sélection
\item Satisfaction accrue des équipes RH
\end{itemize}

\subsubsection{Impact métier attendu}

L'impact attendu sur les processus RH :
\begin{itemize}
\item Réduction des coûts de recrutement de 40-60\%
\item Amélioration de la qualité des embauches
\item Accélération des délais de recrutement
\item Amélioration de l'expérience candidat
\end{itemize}

\subsection{Objectifs de recherche}

Au-delà des objectifs opérationnels, le projet contribue à la recherche en :
\begin{itemize}
\item Adaptation de modèles d'IA à des domaines spécialisés
\item Développement de systèmes hybrides multi-signaux
\item Amélioration de l'explicabilité des systèmes d'IA
\item Évaluation de l'impact des biais algorithmiques
\end{itemize}

\section{Solution proposée}

\subsection{Architecture hybride}

Notre approche combine trois signaux complémentaires :
\begin{enumerate}
\item \textbf{Matching sémantique :} Utilisation d'embeddings denses adaptés au domaine
\item \textbf{Connaissances expertes :} Graphe de proximité des compétences
\item \textbf{Signaux externes :} Extraction d'informations depuis les profils publics
\end{enumerate}

\subsubsection{Architecture détaillée}

\begin{figure}[H]
\centering
\includegraphics[width=0.9\textwidth]{architecture_detaillee.png}
\caption{Architecture hybride du système de matching}
\label{fig:architecture_hybride}
\end{figure}

\subsection{Technologies clés}

Le système repose sur :
\begin{itemize}
\item Modèles de transformers fine-tunés (BAAI BGE)
\item Graphes de connaissances pour les relations compétence-compétence
\item Techniques de web scraping éthique pour l'extraction URL
\item APIs REST pour l'intégration système
\end{itemize}

\subsubsection{Choix technologiques}

Les technologies sélectionnées répondent aux critères suivants :
\begin{itemize}
\item \textbf{Maturité :} Technologies éprouvées et stables
\item \textbf{Performance :} Capacité à traiter de gros volumes
\item \textbf{Ouverture :} Solutions open-source préférées
\item \textbf{Évolutivité :} Architecture cloud-native
\end{itemize}

\subsection{Avantages différenciants}

Par rapport aux solutions existantes, notre approche offre :
\begin{itemize}
\item Une adaptation spécifique au domaine du recrutement
\item Une explicabilité complète et structurée
\item Un support multi-langues natif
\item Une architecture modulaire et extensible
\end{itemize}

\subsubsection{Différenciation concurrentielle}

Les facteurs de différenciation clés :
\begin{itemize}
\item \textbf{Approche hybride :} Combinaison unique de signaux multiples
\item \textbf{Explicabilité native :} Transparence intégrée dès la conception
\item \textbf{Adaptation domaine :} Fine-tuning spécifique au recrutement
\item \textbf{Éthique intégrée :} Considérations RGPD et équité natives
\end{itemize}

\subsection{Stratégie de validation}

La validation du système s'appuie sur :
\begin{itemize}
\item Métriques quantitatives (corrélation, précision, rappel)
\item Évaluation qualitative (explicabilité, utilisabilité)
\item Tests utilisateurs pilotes
\item Comparaison avec baselines établies
\end{itemize}

\subsubsection{Protocole d'évaluation}

Le protocole d'évaluation comprend :
\begin{enumerate}
\item Tests unitaires et d'intégration
\item Validation sur datasets de référence
\item Tests de performance et robustesse
\item Évaluation utilisateur qualitative
\end{enumerate}

% Chapitre 2
\part{Chapitre 2 : Étude de marché et faisabilité}
\label{part:marche}
\chapter{Étude de marché et faisabilité}
\label{chap:marche}

\section{Analyse PESTEL}

\subsection{Politique}

\subsubsection{Réglementation du travail et de l'emploi}

Le secteur du recrutement est fortement régulé par :
\begin{itemize}
\item Les lois sur l'égalité des chances et la non-discrimination (directive européenne 2000/78/CE)
\item Les réglementations sur la protection des données personnelles (RGPD en Europe)
\item Les politiques d'immigration et de mobilité internationale
\item Les normes sur la santé et sécurité au travail
\end{itemize}

\subsubsection{Impact des politiques publiques}

Les initiatives gouvernementales favorisent l'innovation dans le recrutement :
\begin{itemize}
\item Programmes de numérisation des services publics (France 2030)
\item Subventions pour l'adoption de technologies RH
\item Politiques de promotion de l'IA responsable (AI Act européen)
\item Initiatives pour l'inclusion et la diversité
\end{itemize}

\subsubsection{Évolution réglementaire}

L'évolution récente de la régulation inclut :
\begin{itemize}
\item Renforcement des obligations d'explicabilité des algorithmes
\item Interdiction des décisions entièrement automatisées dans certains domaines
\item Obligations de transparence sur les données utilisées
\item Sanctions accrues pour non-conformité
\end{itemize}

\subsubsection{Impact sur notre solution}

Ces évolutions réglementaires influencent notre approche :
\begin{itemize}
\item Intégration native de mécanismes d'explicabilité
\item Conformité RGPD dès la conception
\item Évaluation d'impact algorithmique
\item Mécanismes d'audit et de traçabilité
\end{itemize}

\subsection{Économique}

\subsubsection{Dynamique du marché du recrutement}

Le marché mondial des solutions de recrutement automatisé représente :
\begin{itemize}
\item Un volume de 15 milliards de dollars en 2024
\item Une croissance annuelle de 8-10\% attendue d'ici 2030
\item Une concentration sur les solutions SaaS et cloud-native
\item Une fragmentation avec de nombreux acteurs spécialisés
\end{itemize}

\subsubsection{Facteurs de croissance}

La demande est tirée par :
\begin{itemize}
\item L'accélération de la transformation digitale des entreprises
\item La pénurie de talents dans les secteurs technologiques
\item La pression sur les coûts opérationnels RH (15-20\% des coûts entreprise)
\item L'augmentation du volume de candidatures par poste (200-500 en moyenne)
\end{itemize}

\subsubsection{Modèle économique}

Notre solution s'inscrit dans un modèle SaaS avec :
\begin{itemize}
\item Tarification par volume de matching (0,10€ par CV traité)
\item Offre freemium pour les PME (jusqu'à 100 matchings/mois gratuits)
\item Services premium pour les grandes entreprises
\item Modules complémentaires (analytics, intégrations)
\end{itemize}

\subsubsection{Projection de marché}

L'analyse de marché révèle :
\begin{itemize}
\item Marché adressable : 2,5 milliards d'euros en Europe
\item Taux de pénétration actuel : < 5\% des entreprises
\item Potentiel de croissance : 25\% CAGR sur 5 ans
\item Concurrence : modérée avec barrières techniques élevées
\end{itemize}

\subsection{Socioculturel}

\subsubsection{Évolution des attentes des candidats}

Les candidats modernes privilégient :
\begin{itemize}
\item Des processus de recrutement rapides et transparents
\item Des retours d'information constructifs et personnalisés
\item Une communication digitale fluide et moderne
\item Des expériences utilisateur comparables aux standards tech
\end{itemize}

\subsubsection{Culture d'entreprise et diversité}

Les organisations accordent une importance croissante à :
\begin{itemize}
\item La promotion de la diversité et de l'inclusion (objectifs DEIB)
\item La réduction des biais algorithmiques et cognitifs
\item L'éthique dans l'utilisation des données personnelles
\item La responsabilité sociale des entreprises
\end{itemize}

\subsubsection{Adoption des technologies}

La société civile montre une acceptance croissante de l'IA dans le recrutement, sous réserve de :
\begin{itemize}
\item Transparence des algorithmes utilisés
\item Possibilité de contestation des décisions
\item Protection des droits individuels
\item Démonstration de valeur ajoutée
\end{itemize}

\subsubsection{Impact sociétal}

Notre solution contribue positivement à :
\begin{itemize}
\item Réduction des discriminations dans le recrutement
\item Accélération de l'accès à l'emploi
\item Amélioration de l'équité procédurale
\item Modernisation des pratiques RH
\end{itemize}

\subsection{Technologique}

\subsubsection{État de l'art en IA et NLP}

Les avancées technologiques clés incluent :
\begin{itemize}
\item Les modèles de transformers de grande taille (GPT, BERT, BGE)
\item Les techniques de fine-tuning pour l'adaptation domaine
\item Les approches d'explicabilité (XAI, LIME, SHAP)
\item Les architectures multi-modales (texte + données structurées)
\end{itemize}

\subsubsection{Infrastructure cloud}

L'écosystème technique bénéficie de :
\begin{itemize}
\item Services de calcul élastique (AWS, GCP, Azure)
\item APIs d'IA managées et accessibles
\item Outils de MLOps pour le déploiement continu
\item Solutions de stockage et traitement distribué
\end{itemize}

\subsubsection{Tendances émergentes}

Les innovations à surveiller incluent :
\begin{itemize}
\item L'IA générative pour la génération d'explications
\item Les modèles multi-modaux (texte + image + CV visuel)
\item L'apprentissage continu et l'adaptation en ligne
\item Les techniques de federated learning pour la confidentialité
\end{itemize}

\subsubsection{Positionnement technologique}

Notre solution s'appuie sur :
\begin{itemize}
\item Technologies matures et éprouvées
\item Stack technique moderne et maintenable
\item Approche hybride combinant plusieurs paradigmes
\item Architecture cloud-native et scalable
\end{itemize}

\subsection{Environnemental}

\subsubsection{Impact environnemental des solutions cloud}

L'adoption de solutions cloud implique :
\begin{itemize}
\item Une consommation énergétique des data centers
\item Des émissions de CO2 liées au calcul intensif
\item Des stratégies de compensation carbone
\item Une optimisation des ressources computationnelles
\end{itemize}

\subsubsection{Durabilité des pratiques}

Les solutions durables doivent considérer :
\begin{itemize}
\item L'optimisation des ressources computationnelles
\item La réduction de la latence réseau
\item La minimisation du stockage redondant
\item L'efficacité énergétique des algorithmes
\end{itemize}

\subsubsection{Contribution environnementale}

Notre approche contribue à la durabilité par :
\begin{itemize}
\item Réduction du papier et des impressions physiques
\item Optimisation des processus RH (moins de déplacements)
\item Utilisation d'infrastructures cloud efficientes
\item Mesure et réduction de l'empreinte carbone
\end{itemize}

\subsection{Légal}

\subsubsection{Conformité RGPD et protection des données}

Les contraintes réglementaires majeures :
\begin{itemize}
\item Consentement explicite pour le traitement des données
\item Droit à l'explication des décisions algorithmiques
\item Obligation de transparence sur les sources de données
\item Droit à la portabilité et à l'effacement des données
\end{itemize}

\subsubsection{Réglementation sur l'IA}

Le cadre européen sur l'IA classe les systèmes de recrutement comme :
\begin{itemize}
\item Systèmes à haut risque nécessitant évaluation
\item Obligation d'audit et de documentation
\item Interdiction des décisions entièrement automatisées dans certains cas
\item Exigences de transparence et d'explicabilité
\end{itemize}

\subsubsection{Propriété intellectuelle}

Les aspects juridiques à considérer :
\begin{itemize}
\item Protection des modèles d'IA développés
\item Licences des datasets et modèles open-source utilisés
\item Droits d'auteur sur les explications générées
\item Brevets sur les algorithmes innovants
\end{itemize}

\subsubsection{Conformité internationale}

Pour un déploiement international :
\begin{itemize}
\item Conformité CCPA (Californie) pour les données US
\item Respect des réglementations locales (ex: PDPA Singapour)
\item Adaptation aux standards sectoriels (ex: SOX pour la finance)
\end{itemize}

\section{Analyse SWOT}

\subsection{Forces}

\subsubsection{Avantages technologiques}

\begin{itemize}
\item \textbf{Précision supérieure :} Corrélation Spearman de 0,78 vs 0,62 pour les baselines
\item \textbf{Explicabilité native :} Justifications structurées pour chaque décision
\item \textbf{Adaptabilité domaine :} Fine-tuning spécifique au recrutement
\item \textbf{Support multi-langues :} Couverture de 100+ langues via BGE-M3
\item \textbf{Performance :} Traitement de 100 CV en < 30 secondes
\end{itemize}

\subsubsection{Avantages opérationnels}

\begin{itemize}
\item \textbf{Évolutivité :} Architecture cloud-native et modulaire
\item \textbf{Intégration :} API REST standard et documentation OpenAPI
\item \textbf{Conformité :} RGPD et éthique intégrés dès la conception
\item \textbf{Support :} Documentation complète et communauté active
\end{itemize}

\subsubsection{Avantages concurrentiels}

\begin{itemize}
\item \textbf{Différenciation :} Approche hybride unique multi-signaux
\item \textbf{Innovation :} Combinaison de techniques de pointe
\item \textbf{Reproductibilité :} Code open-source et méthodologie documentée
\end{itemize}

\subsection{Faiblesses}

\subsubsection{Limitations techniques}

\begin{itemize}
\item \textbf{Dépendance aux données :} Performance liée à la qualité des données d'entraînement
\item \textbf{Complexité d'implémentation :} Expertise requise en ML et NLP
\item \textbf{Coûts computationnels :} Ressources nécessaires pour le fine-tuning
\item \textbf{Maintenance :} Nécessité de mises à jour régulières des modèles
\end{itemize}

\subsubsection{Contraintes opérationnelles}

\begin{itemize}
\item \textbf{Adoption culturelle :} Résistance au changement dans les équipes RH
\item \textbf{Courbe d'apprentissage :} Formation nécessaire pour l'utilisation optimale
\item \textbf{Dépendance fournisseur :} Liens avec les providers cloud et de données
\end{itemize}

\subsubsection{Risques identifiés}

\begin{itemize}
\item \textbf{Obsolescence :} Évolution rapide des technologies d'IA
\item \textbf{Concurrence :} Arrivée de nouveaux acteurs mieux financés
\item \textbf{Réglementation :} Changements dans les lois sur l'IA
\end{itemize}

\subsection{Oportunités}

\subsubsection{Marché porteur}

\begin{itemize}
\item \textbf{Croissance du secteur :} Marché de l'IA RH en expansion rapide (+25\% CAGR)
\item \textbf{Demande institutionnelle :} Intérêt croissant des grandes entreprises
\item \textbf{Réglementation favorable :} Incitations à l'adoption de solutions éthiques
\item \textbf{Transformation digitale :} Accélération post-COVID des investissements IT
\end{itemize}

\subsubsection{Extensions possibles}

\begin{itemize}
\item \textbf{Diversification sectorielle :} Adaptation à d'autres domaines (santé, finance, éducation)
\item \textbf{Écosystème :} Développement d'APIs et intégrations tierces
\item \textbf{Internationalisation :} Expansion géographique et linguistique
\item \textbf{Innovation :} Recherche de nouveaux cas d'usage
\end{itemize}

\subsubsection{Partenariats stratégiques}

\begin{itemize}
\item \textbf{Éditeurs RH :} Intégration avec Workday, SAP SuccessFactors
\item \textbf{Plateformes :} Partenariats avec LinkedIn, Indeed
\item \textbf{Universités :} Collaborations pour la recherche et les datasets
\item \textbf{Investisseurs :} Levée de fonds pour l'accélération
\end{itemize}

\subsection{Menaces}

\subsubsection{Concurrence}

\begin{itemize}
\item \textbf{Acteurs établis :} Solutions propriétaires des grands éditeurs (Workday, SAP)
\item \textbf{Startups innovantes :} Nouvelles entreprises spécialisées en IA RH
\item \textbf{Solutions open-source :} Alternatives gratuites et communautaires
\item \textbf{Internalisation :} Développement interne par les grandes entreprises
\end{itemize}

\subsubsection{Risques technologiques}

\begin{itemize}
\item \textbf{Évolution rapide :} Obsolescence des modèles face aux avancées (GPT-5, etc.)
\item \textbf{Dépendance fournisseur :} Risques liés aux APIs externes (GitHub, LinkedIn)
\item \textbf{Failles de sécurité :} Vulnérabilités dans les chaînes d'approvisionnement
\item \textbf{Attaques adversarielles :} Manipulation des inputs pour tromper le système
\end{itemize}

\subsubsection{Risques réglementaires}

\begin{itemize}
\item \textbf{Évolution légale :} Changements dans la régulation de l'IA
\item \textbf{Contentieux :} Actions en justice liées à des décisions discriminatoires
\item \textbf{Amendements :} Modifications des lois sur la protection des données
\end{itemize}

\subsubsection{Risques économiques}

\begin{itemize}
\item \textbf{Récession :} Réduction des budgets RH en période de crise
\item \textbf{Concurrence prix :} Pression sur les marges par les solutions low-cost
\item \textbf{Change rate :} Impact des fluctuations monétaires sur l'international
\end{itemize}

\section{Analyse des 5 forces de Porter}

\subsection{Concurrence}

\subsubsection{Structure concurrentielle}

Le marché présente une concurrence modérée caractérisée par :
\begin{itemize}
\item Quelques acteurs dominants (LinkedIn, Indeed, Glassdoor)
\item Nombreuses startups spécialisées en IA RH
\item Solutions open-source émergentes
\item Développement interne par les grandes entreprises
\end{itemize}

\subsubsection{Positionnement des concurrents}

\begin{table}[H]
\centering
\caption{Analyse comparative des concurrents}
\label{tab:concurrents}
\begin{tabular}{@{}lllll@{}}
\toprule
Acteur & Modèle & Forces & Faiblesses & Positionnement \\
\midrule
LinkedIn & Plateforme & Réseau étendu & Coûts élevés & Dominant \\
Indeed & Job board & Volume & Matching basique & Volume \\
Workday & Suite RH & Intégration & Complexité & Enterprise \\
Startups IA & Spécialisées & Innovation & Ressources & Niche \\
Notre solution & SaaS IA & Explicabilité & Adoption & Différencié \\
\bottomrule
\end{tabular}
\end{table}

\subsubsection{Différenciation}

Notre positionnement se distingue par :
\begin{itemize}
\item Focus exclusif sur l'explicabilité et la transparence
\item Approche hybride combinant plusieurs signaux
\item Support multi-langues natif
\item Architecture ouverte et extensible
\end{itemize}

\subsubsection{Barrières à l'entrée}

Les obstacles pour les nouveaux entrants :
\begin{itemize}
\item Expertise technique spécialisée (ML/NLP)
\item Accès à des datasets de qualité
\item Conformité réglementaire complexe
\item Réseau de distribution établi
\end{itemize}

\subsection{Pouvoir des clients}

\subsubsection{Caractéristiques des acheteurs}

Les clients (recruteurs, RH, managers) exercent un pouvoir significatif :
\begin{itemize}
\item \textbf{Concentration :} Grandes entreprises représentent la majorité du marché
\item \textbf{Information :} Accès facile aux alternatives et comparaisons
\item \textbf{Sensibilité prix :} Budgets RH sous pression économique
\item \textbf{Exigences élevées :} Attentes de performance et de qualité
\end{itemize}

\subsubsection{ Segmentation des clients}

\begin{itemize}
\item \textbf{Grandes entreprises (>1000 employés) :} 40\% du marché, besoins complexes
\item \textbf{ETI (200-1000 employés) :} 35\%, recherche de solutions intégrées
\item \textbf{PME (20-200 employés) :} 20\%, sensibilité prix, besoins simples
\item \textbf{Startups (<20 employés) :} 5\%, innovation, contraintes budgétaires
\end{itemize}

\subsubsection{Leverage client}

Les stratégies pour renforcer notre position :
\begin{itemize}
\item Offre freemium pour réduire les barrières d'entrée
\item Démonstrations personnalisées et POC
\item Support client premium et formation
\item Intégrations natives avec les outils RH existants
\end{itemize}

\subsubsection{Négociation client}

Le pouvoir de négociation des clients dépend de :
\begin{itemize}
\item Taille de l'entreprise (plus grande = plus de pouvoir)
\item Nombre d'alternatives disponibles
\item Coûts de changement (switching costs)
\item Degré de différenciation perçu
\end{itemize}

\subsection{Pouvoir des fournisseurs}

\subsubsection{Dépendance technologique}

Notre dépendance concerne principalement :
\begin{itemize}
\item Fournisseurs de modèles d'IA (HuggingFace, OpenAI)
\item Services cloud (AWS, GCP)
\item APIs externes (GitHub, LinkedIn)
\item Datasets et bibliothèques open-source
\end{itemize}

\subsubsection{Évaluation des dépendances}

\begin{table}[H]
\centering
\caption{Évaluation des dépendances fournisseurs}
\label{tab:fournisseurs}
\begin{tabular}{@{}llll@{}}
\toprule
Fournisseur & Produit & Niveau dépendance & Stratégie mitigation \\
\midrule
HuggingFace & Modèles & Élevé & Multi-fournisseurs \\
AWS/GCP & Cloud & Moyen & Migration possible \\
GitHub & API & Faible & Alternatives (GitLab) \\
Open-source & Bibliothèques & Faible & Forks disponibles \\
\bottomrule
\end{tabular}
\end{table}

\subsubsection{Stratégies de mitigation}

Pour réduire la dépendance fournisseur :
\begin{itemize}
\item Diversification des sources de modèles
\item Développement de capacités internes
\item Stockage de modèles locaux et cache
\item Contrats de service garantis
\end{itemize}

\subsubsection{Négociation fournisseur}

Notre pouvoir de négociation repose sur :
\begin{itemize}
\item Volume d'utilisation potentiel
\item Contribution à la communauté open-source
\item Relations établies et références
\item Alternatives techniques disponibles
\end{itemize}

\subsection{Menace des nouveaux entrants}

\subsubsection{Barrières à l'entrée}

Les obstacles pour de nouveaux entrants sont élevés :
\begin{itemize}
\item \textbf{Expertise technique :} Compétences avancées en ML/NLP requises
\item \textbf{Données d'entraînement :} Accès à des datasets de qualité
\item \textbf{Réglementation :} Conformité RGPD et normes IA
\item \textbf{Capitaux :} Investissements nécessaires pour le développement
\item \textbf{Réseau :} Relations établies avec les clients RH
\end{itemize}

\subsubsection{Évaluation des barrières}

\begin{table}[H]
\centering
\caption{Évaluation des barrières à l'entrée}
\label{tab:barrieres}
\begin{tabular}{@{}lll@{}}
\toprule
Barrière & Hauteur & Commentaire \\
\midrule
Technique & Très élevée & Expertise ML/NLP spécialisée \\
Données & Élevée & Datasets annotés coûteux \\
Réglementaire & Élevée & Conformité RGPD/AI Act \\
Financière & Moyenne & Coûts développement 500K€ \\
Marketing & Moyenne & Réseau RH à construire \\
\bottomrule
\end{tabular}
\end{table}

\subsubsection{Vulnérabilités}

Malgré ces barrières, les menaces incluent :
\begin{itemize}
\item Startups bien financées avec expertise pointue
\item Grandes tech companies entrant sur le marché RH
\item Solutions open-source matures
\item Développement interne par les entreprises
\end{itemize}

\subsubsection{Évolution des barrières}

Les barrières pourraient s'abaisser avec :
\begin{itemize}
\item Démocratisation des outils ML (AutoML)
\item Disponibilité de datasets open-source
\item Simplification des réglementations
\item Réduction des coûts de calcul
\end{itemize}

\subsection{Produits de substitution}

\subsubsection{Alternatives disponibles}

Les substituts incluent :
\begin{itemize}
\item \textbf{Méthodes manuelles :} Screening traditionnel par les recruteurs
\item \textbf{ATS classiques :} Solutions sans IA avancée (Bullhorn, Taleo)
\item \textbf{Plateformes généralistes :} LinkedIn, Indeed avec matching basique
\item \textbf{Solutions internes :} Développement propriétaire par les entreprises
\end{itemize}

\subsubsection{Analyse des substituts}

\begin{table}[H]
\centering
\caption{Analyse des produits de substitution}
\label{tab:substituts}
\begin{tabular}{@{}llll@{}}
\toprule
Substitut & Avantages & Inconvénients & Adoption \\
\midrule
Manuel & Contrôle total & Coûteux, lent & Décroissant \\
ATS classique & Établi, intégré & Matching limité & Stable \\
Plateformes & Facile d'accès & Générique & Croissant \\
Interne & Personnalisé & Coûteux à développer & Niche \\
\bottomrule
\end{tabular}
\end{table}

\subsubsection{Avantage concurrentiel}

Notre différenciation repose sur :
\begin{itemize}
\item Précision supérieure démontrée
\item Explicabilité complète
\item Automatisation sans perte de contrôle humain
\item Coûts compétitifs pour la valeur apportée
\end{itemize}

\subsubsection{Évolution des substitutions}

Les tendances qui affectent les substituts :
\begin{itemize}
\item Augmentation des volumes de candidatures
\item Pression sur les coûts RH
\item Demande d'équité et de transparence
\item Adoption croissante de l'IA
\end{itemize}

\section{Étude de faisabilité}

\subsection{Étude technique}

\subsubsection{Faisabilité technologique}

L'analyse technique confirme la viabilité du projet :
\begin{itemize}
\item Technologies matures et accessibles (transformers, embeddings)
\item Infrastructure cloud disponible et scalable
\item Expertise disponible sur le marché
\item Outils de développement adaptés
\end{itemize}

\subsubsection{Risques techniques identifiés}

\begin{itemize}
\item Complexité du fine-tuning sur données limitées
\item Performance des modèles multi-langues
\item Robustesse de l'extraction URL
\item Gestion des biais algorithmiques
\end{itemize}

\subsubsection{Plan de mitigation technique}

\begin{itemize}
\item Validation expérimentale précoce des approches
\item Benchmarks comparatifs systématiques
\item Prototypage itératif et tests pilotes
\item Revue de code et bonnes pratiques
\end{itemize}

\subsubsection{Ressources techniques requises}

\begin{itemize}
\item Équipe de développement (4-6 personnes)
\item Infrastructure cloud (budget 50K€/an)
\item Datasets et modèles (licences open-source)
\item Outils de développement et CI/CD
\end{itemize}

\subsection{Étude commerciale}

\subsubsection{Marché cible}

Le marché adressable comprend :
\begin{itemize}
\item Grandes entreprises (>1000 employés) : 40\% du marché
\item PME technologiques (100-1000 employés) : 35\%
\item Organisations publiques et éducation : 25\%
\end{itemize}

\subsubsection{Segmentation détaillée}

\begin{table}[H]
\centering
\caption{Segmentation du marché cible}
\label{tab:segmentation}
\begin{tabular}{@{}llll@{}}
\toprule
Segment & Taille & Caractéristiques & Offre adaptée \\
\midrule
Grandes entreprises & 40\% & Besoins complexes, budget élevé & Solution enterprise complète \\
PME tech & 35\% & Innovation, contraintes budget & SaaS flexible, freemium \\
Secteur public & 25\% & Réglementation stricte, transparence & Solution certifiée, explicable \\
\bottomrule
\end{tabular}
\end{table}

\subsubsection{Modèle de revenus}

Stratégie de monétisation :
\begin{itemize}
\item Abonnement SaaS : 70\% des revenus (100-500€/mois selon taille)
\item Services professionnels : 20\% (consulting, intégration)
\item Licences enterprise : 10\% (déploiement on-premise)
\end{itemize}

\subsubsection{Projection financière}

Estimations sur 3 ans :
\begin{itemize}
\item CA année 1 : 500K€
\item CA année 2 : 2M€
\item CA année 3 : 5M€
\item Point mort : 18 mois
\end{itemize}

\subsubsection{Stratégie de go-to-market}

\begin{itemize}
\item Phase 1 : Validation produit avec early adopters
\item Phase 2 : Expansion via partenariats et marketing digital
\item Phase 3 : Vente enterprise et internationalisation
\end{itemize}

\subsection{Étude financière}

\subsubsection{Investissement initial}

Budget de développement estimé :
\begin{itemize}
\item Ressources humaines : 400K€ (salaires équipe 24 mois)
\item Infrastructure cloud : 50K€
\item Données et datasets : 25K€
\item Marketing et commercial : 75K€
\item Total : 550K€
\end{itemize}

\subsubsection{Charges opérationnelles}

Coûts récurrents annuels :
\begin{itemize}
\item Infrastructure : 100K€
\item Support et maintenance : 150K€
\item R\&D continue : 200K€
\item Marketing : 150K€
\item Total : 600K€
\end{itemize}

\subsubsection{Analyse de rentabilité}

\begin{itemize}
\item Chiffre d'affaires cible année 1 : 500K€
\item Charges année 1 : 600K€
\item Résultat année 1 : -100K€
\item ROI positif dès l'année 2
\item Marges opérationnelles : 60-70\% en régime de croisière
\end{itemize}

\subsubsection{Risques financiers}

\begin{itemize}
\item Dépendance à l'adoption client (risque principal)
\item Évolution des coûts cloud
\item Concurrence prix des solutions alternatives
\item Retards de développement
\end{itemize}

\subsubsection{Plan de financement}

\begin{itemize}
\item Apport personnel : 100K€
\item Levée de fonds : 300K€ (amorçage)
\item Subventions : 100K€ (projets IA européens)
\item Crédits : 50K€ (prêts innovation)
\end{itemize}

\subsection{Étude opérationnelle}

\subsubsection{Plan de développement}

Chronologie sur 24 mois :
\begin{enumerate}
\item Mois 1-6 : Développement core et validation
\item Mois 7-12 : Tests pilotes et itérations
\item Mois 13-18 : Lancement commercial et expansion
\item Mois 19-24 : Optimisation et internationalisation
\end{enumerate}

\subsubsection{Ressources humaines}

Équipe cible :
\begin{itemize}
\item Lead ML Engineer : Architecture IA
\item Backend Developer : APIs et infrastructure
\item Frontend Developer : Interface utilisateur
\item Product Manager : Stratégie et roadmap
\item DevOps Engineer : Déploiement et monitoring
\end{itemize}

\subsubsection{Risques opérationnels}

\begin{itemize}
\item Recrutement de talents spécialisés
\item Gestion de la croissance technique
\item Maintenance de la qualité du code
\item Gestion des relations clients
\end{itemize}

\subsection{Conclusion de faisabilité}

\subsubsection{Synthèse}

L'étude de faisabilité conclut à la viabilité du projet :
\begin{itemize}
\item \textbf{Technique :} Faisable avec les technologies actuelles
\item \textbf{Commerciale :} Marché porteur avec demande croissante
\item \textbf{Financière :} Modèle rentable avec ROI positif à 2 ans
\item \textbf{Operationnelle :} Équipe et processus réalisables
\end{itemize}

\subsubsection{Conditions de succès}

\begin{itemize}
\item Validation technique des performances promises
\item Adoption par les premiers clients pilotes
\item Sécurisation du financement initial
\item Constitution d'une équipe compétente
\end{itemize}

\subsubsection{Recommandations}

\begin{itemize}
\item Démarrer par un MVP focalisé sur les fonctionnalités core
\item Valider le produit-market fit avec 3-5 clients pilotes
\item Préparer une stratégie de scaling précoce
\item Maintenir une veille technologique active
\end{itemize}

% Chapitre 3
\part{Chapitre 3 : État de l'art et revue de littérature}
\label{part:etat-art}
\chapter{État de l'art et revue de littérature}
\label{chap:etat-art}

\section{Évolution historique du matching CV}

\subsection{Les approches traditionnelles}

\subsubsection{Matching par mots-clés (avant 2010)}

Les premières approches de matching CV reposaient sur des techniques simples :
\begin{itemize}
\item Recherche exacte de termes dans les CV et descriptions de poste
\item Score TF-IDF pour mesurer la fréquence des termes
\item Matching booléen (présence/absence de mots-clés)
\item Systèmes de règles manuelles définies par les recruteurs
\end{itemize}

\subsubsection{Limitations des approches classiques}

Ces méthodes présentaient plusieurs contraintes majeures :
\begin{itemize}
\item \textbf{Faible précision :} Ignorent le contexte sémantique
\item \textbf{Sensibilité aux variations :} "développeur" ≠ "développement"
\item \textbf{Biais linguistiques :} Dépendent de la formulation exacte
\item \textbf{Manque d'adaptabilité :} Règles fixes non évolutives
\end{itemize}

\subsubsection{Impact sur l'efficacité recrutement}

Selon les études de l'époque :
\begin{itemize}
\item Taux de rejet des bons candidats : 20-30\%
\item Temps moyen de screening : 15-20 minutes par CV
\item Satisfaction des recruteurs : < 60\%
\item Coûts de recrutement : 15-25\% du salaire annuel
\end{itemize}

\subsection{L'avènement du machine learning (2010-2015)}

\subsubsection{Approches supervisées}

L'introduction du machine learning a apporté :
\begin{itemize}
\item Classifieurs SVM pour la prédiction de compatibilité
\item Réseaux de neurones pour l'extraction de features
\item Apprentissage supervisé sur données annotées
\item Métriques d'évaluation standardisées (précision, rappel, F1-score)
\end{itemize}

\subsubsection{Extraction de features}

Les techniques d'extraction incluaient :
\begin{itemize}
\item Bag-of-words avec pondération TF-IDF
\item N-grams pour capturer les expressions
\item Features syntaxiques (POS tagging)
\item Features sémantiques basiques
\end{itemize}

\subsubsection{Résultats obtenus}

Les performances s'améliorent significativement :
\begin{itemize}
\item Précision : 65-75\% vs 50-60\% pour les approches classiques
\item Rappel : 70-80\% avec réduction des faux négatifs
\item F1-score : 0,72 en moyenne
\item Réduction du temps de screening : 40-50\%
\end{itemize}

\subsection{La révolution des embeddings (2015-2020)}

\subsubsection{Word embeddings}

L'arrivée des word embeddings transforme l'approche :
\begin{itemize}
\item Word2Vec (2013) : représentations vectorielles denses
\item GloVe (2014) : apprentissage sur matrices de co-occurrence
\item FastText (2016) : gestion des mots hors-vocabulaire
\item Embeddings contextuels (ELMo, 2018)
\end{itemize}

\subsubsection{Applications au matching CV}

Les embeddings permettent :
\begin{itemize}
\item Similarité sémantique entre CV et offres
\item Représentation de compétences équivalentes
\item Gestion des synonymes et paraphrases
\item Réduction de la dimensionnalité
\end{itemize}

\subsubsection{Performances améliorées}

Les résultats montrent des gains significatifs :
\begin{itemize}
\item Corrélation avec jugements humains : 0,65-0,75
\item Précision sémantique : +15-20\% vs approches précédentes
\item Robustesse aux reformulations : +30\%
\item Gestion multi-langues : premières approches
\end{itemize}

\subsection{L'ère des transformers (2020-présent)}

\subsubsection{Architecture transformer}

Les transformers révolutionnent le domaine :
\begin{itemize}
\item Attention mechanism (Vaswani et al., 2017)
\item BERT (2018) : bidirectional encoding
\item GPT (2018) : génération autoregressive
\item BART et T5 (2020) : architectures seq2seq
\end{itemize}

\subsubsection{Applications spécifiques au recrutement}

Dans le domaine RH, les transformers permettent :
\begin{itemize}
\item Compréhension contextuelle des descriptions
\item Extraction fine des compétences et expériences
\item Génération d'explications justifiées
\item Adaptation domaine-spécifique via fine-tuning
\end{itemize}

\subsubsection{État actuel des performances}

Les systèmes modernes atteignent :
\begin{itemize}
\item Corrélation Spearman : 0,75-0,85
\item Précision métier : 80-90\% sur critères pertinents
\item Explicabilité : justifications textuelles cohérentes
\item Scalabilité : traitement de milliers de CV/minute
\end{itemize}

\section{Techniques de base en NLP pour le matching}

\subsection{Prétraitement du texte}

\subsubsection{Normalisation}

Les étapes de normalisation incluent :
\begin{itemize}
\item Conversion en minuscules
\item Suppression de la ponctuation et caractères spéciaux
\item Normalisation Unicode (décomposition canonique)
\item Gestion des encodages multiples
\end{itemize}

\subsubsection{Tokenisation}

La tokenisation avancée comprend :
\begin{itemize}
\item Segmentation en mots et sous-mots
\item Gestion des contractions ("j'ai" → "je ai")
\item Traitement des entités nommées (noms, dates)
\item Conservation des informations de casse pour NER
\end{itemize}

\subsubsection{Filtrage et nettoyage}

Le nettoyage des données implique :
\begin{itemize}
\item Suppression des stop-words (articles, prépositions)
\item Filtrage des termes trop fréquents ou rares
\item Correction orthographique automatique
\item Détection et suppression du bruit (headers, footers)
\end{itemize}

\subsection{Extraction d'informations structurées}

\subsubsection{Extraction de compétences}

Les techniques d'extraction incluent :
\begin{itemize}
\item Pattern matching avec expressions régulières
\item Modèles CRF (Conditional Random Fields)
\item Réseaux de neurones pour NER (Named Entity Recognition)
\item Ontologies de compétences (ESCO, O*NET)
\end{itemize}

\subsubsection{Extraction d'expériences}

Pour les expériences professionnelles :
\begin{itemize}
\item Analyse syntaxique pour identifier les structures
\item Extraction de dates et durées
\item Identification des rôles et responsabilités
\item Classification des niveaux d'expérience
\end{itemize}

\subsubsection{Extraction d'éducation}

L'extraction des formations comprend :
\begin{itemize}
\item Reconnaissance des diplômes et certifications
\item Classification des niveaux d'études
\item Extraction des domaines d'études
\item Validation des établissements
\end{itemize}

\subsection{Représentation vectorielle}

\subsubsection{Approches traditionnelles}

Les méthodes classiques incluent :
\begin{itemize}
\item Bag-of-Words (BoW) avec vocabulaire fixe
\item TF-IDF pour pondération des termes importants
\item N-grams pour capturer les expressions locales
\item Hashing trick pour réduction dimensionnelle
\end{itemize}

\subsubsection{Embeddings statiques}

Les word embeddings statiques :
\begin{itemize}
\item Word2Vec : CBOW et Skip-gram architectures
\item GloVe : factorisation de matrices de co-occurrence
\item FastText : incorporation des sous-mots
\item Avantages : légèreté et rapidité d'inférence
\end{itemize}

\subsubsection{Embeddings contextuels}

Les représentations contextuelles modernes :
\begin{itemize}
\item ELMo : concaténation de couches bidirectionnelles
\item BERT : transformers bidirectionnels pré-entraînés
\item RoBERTa : optimisation de BERT
\item DistilBERT : version distillée pour l'efficacité
\end{itemize}

\subsubsection{Embeddings spécialisés recrutement}

Pour le domaine RH spécifiquement :
\begin{itemize}
\item SciBERT : adapté aux textes scientifiques
\item BioBERT : spécialisé en biomédical
\item JobBERT : fine-tuné sur données recrutement
\item Domain adaptation via fine-tuning
\end{itemize}

\section{Approches de similarité et matching}

\subsection{Mesures de similarité traditionnelles}

\subsubsection{Distance cosinus}

La similarité cosinus mesure l'angle entre vecteurs :
\begin{equation}
similarity(u,v) = \frac{u \cdot v}{\|u\| \|v\|}
\end{equation}

Avantages :
\begin{itemize}
\item Invariant à la magnitude des vecteurs
\item Interprétable géométriquement
\item Efficace computationnellement
\end{itemize}

\subsubsection{Distance euclidienne}

La distance euclidienne :
\begin{equation}
distance(u,v) = \sqrt{\sum_{i=1}^{n}(u_i - v_i)^2}
\end{equation}

Utilisée pour :
\begin{itemize}
\item Mesures absolues de dissimilarité
\item Clustering hiérarchique
\item Détection d'outliers
\end{itemize}

\subsubsection{Similarité Jaccard}

Pour ensembles discrets :
\begin{equation}
J(A,B) = \frac{|A \cap B|}{|A \cup B|}
\end{equation}

Particulièrement utile pour :
\begin{itemize}
\item Comparaison de sets de compétences
\item Matching par mots-clés
\item Évaluation de chevauchement
\end{itemize}

\subsection{Approches par apprentissage automatique}

\subsubsection{Apprentissage supervisé}

Les modèles supervisés incluent :
\begin{itemize}
\item Régression logistique pour scoring binaire
\item SVM avec kernels RBF pour séparation non-linéaire
\item Random Forests pour robustesse aux outliers
\item Gradient Boosting (XGBoost, LightGBM)
\end{itemize}

\subsubsection{Apprentissage non supervisé}

Les méthodes non supervisées :
\begin{itemize}
\item K-means pour clustering de profils
\item PCA pour réduction dimensionnelle
\item Autoencodeurs pour apprentissage de représentations
\item Topic modeling (LDA) pour extraction de thèmes
\end{itemize}

\subsubsection{Apprentissage par renforcement}

Applications émergentes :
\begin{itemize}
\item Optimisation de stratégies de matching
\item Apprentissage de politiques d'interaction
\item Personnalisation adaptative des résultats
\end{itemize}

\subsection{Approches neuronales avancées}

\subsubsection{Réseaux siamois}

Architecture pour apprentissage de similarité :
\begin{itemize}
\item Partage des poids entre deux branches
\item Fonction de perte contrastive
\item Triplet loss pour apprentissage relatif
\item Applications au matching one-shot
\end{itemize}

\subsubsection{Transformers pour matching}

Utilisation des transformers :
\begin{itemize}
\item Cross-encoders pour scoring direct
\item Bi-encoders pour efficacité
\item Late interaction pour précision
\item Hybrid approaches combinant les deux
\end{itemize}

\subsubsection{Approches multi-modales}

Intégration de multiples sources :
\begin{itemize}
\item Texte + métadonnées structurées
\item Incorporation de réseaux sociaux
\item Analyse de vidéos d'entretien
\item Fusion de scores multiples
\end{itemize}

\section{Évaluation et métriques}

\subsection{Métriques classiques}

\subsubsection{Précision et rappel}

Définitions fondamentales :
\begin{equation}
Precision = \frac{TP}{TP + FP}, \quad Recall = \frac{TP}{TP + FN}
\end{equation}

\subsubsection{F1-score}

Moyenne harmonique :
\begin{equation}
F1 = 2 \cdot \frac{Precision \cdot Recall}{Precision + Recall}
\end{equation}

\subsubsection{Matrice de confusion}

Analyse détaillée des erreurs :
\begin{itemize}
\item Vrais positifs (TP) : bons matches identifiés
\item Faux positifs (FP) : mauvais candidats sélectionnés
\item Vrais négatifs (TN) : mauvais candidats rejetés
\item Faux négatifs (FN) : bons candidats manqués
\end{itemize}

\subsection{Métriques métier}

\subsubsection{Corrélation avec jugements humains}

Évaluation par experts :
\begin{itemize}
\item Accord inter-juges (Cohen's Kappa)
\item Corrélation de Spearman avec rankings humains
\item Analyse des désaccords systématiques
\end{itemize}

\subsubsection{Métriques de diversité}

Évaluation de l'équité :
\begin{itemize}
\item Répartition démographique des résultats
\item Analyse des biais algorithmiques
\item Mesure de la diversité des recommandations
\end{itemize}

\subsubsection{Métriques temporelles}

Performance opérationnelle :
\begin{itemize}
\item Temps de traitement par CV
\item Débit maximum du système
\item Latence acceptable pour l'utilisateur
\end{itemize}

\subsection{Protocoles d'évaluation}

\subsubsection{Cross-validation}

Validation croisée pour robustesse :
\begin{itemize}
\item K-fold cross-validation
\item Stratification par domaine
\item Validation temporelle pour éviter le data leakage
\end{itemize}

\subsubsection{Évaluation sur benchmarks}

Benchmarks standardisés :
\begin{itemize}
\item TREC Career Tracking (2009-2011)
\item RESUME (2012)
\item REL+ED (2016)
\item Benchmarks modernes sur données réelles
\end{itemize}

\subsubsection{A/B testing}

Évaluation en production :
\begin{itemize}
\item Comparaison avec système existant
\item Mesures d'engagement utilisateur
\item Impact sur métriques business (time-to-hire, quality-of-hire)
\end{itemize}

\section{Explicabilité et interprétabilité}

\subsection{Importance de l'explicabilité}

\subsubsection{Exigences réglementaires}

L'explicabilité est devenue obligatoire :
\begin{itemize}
\item RGPD : droit à l'explication des décisions
\item AI Act européen : exigences de transparence
\item Lois sectorielles (finance, santé)
\end{itemize}

\subsubsection{Besoins métier}

Les utilisateurs exigent :
\begin{itemize}
\item Compréhension des décisions du système
\item Possibilité de contestation des résultats
\item Confiance dans les recommandations
\item Auditabilité des processus
\end{itemize}

\subsubsection{Avantages techniques}

L'explicabilité apporte :
\begin{itemize}
\item Débogage facilité des modèles
\item Détection des biais et erreurs
\item Amélioration itérative des performances
\item Validation des décisions par domaine expert
\end{itemize}

\subsection{Techniques d'explicabilité}

\subsubsection{Approches globales}

Explication du modèle dans son ensemble :
\begin{itemize}
\item Feature importance (permutation, SHAP)
\item Partial dependence plots
\item Accumulated local effects
\item Global surrogate models
\end{itemize}

\subsubsection{Approches locales}

Explication de prédictions individuelles :
\begin{itemize}
\item LIME (Local Interpretable Model-agnostic Explanations)
\item SHAP values pour attribution locale
\item Anchors pour règles interprétables
\item Counterfactual explanations
\end{itemize}

\subsubsection{Explicabilité pour transformers}

Techniques spécifiques aux transformers :
\begin{itemize}
\item Attention weights visualization
\item Layer-wise relevance propagation
\item Integrated gradients
\item Concept activation vectors
\end{itemize}

\subsection{Génération d'explications}

\subsubsection{Explications textuelles}

Génération automatique :
\begin{itemize}
\item Templates prédéfinis avec variables
\item Modèles de langage pour génération naturelle
\item Extraction de passages pertinents
\item Résumé automatique des raisons
\end{itemize}

\subsubsection{Explications visuelles}

Représentations graphiques :
\begin{itemize}
\item Heatmaps de similarité
\item Graphiques de contribution des features
\item Visualisation des embeddings (t-SNE, UMAP)
\item Decision trees interprétables
\end{itemize}

\subsubsection{Explications interactives}

Interfaces utilisateur :
\begin{itemize}
\item Exploration interactive des résultats
\item "What-if" analysis pour scénarios alternatifs
\item Feedback utilisateur intégré
\item Audit trails complets
\end{itemize}

\section{Défis et limitations actuelles}

\subsection{Défis techniques}

\subsubsection{Gestion des données}

Les contraintes majeures :
\begin{itemize}
\item Disponibilité de données annotées de qualité
\item Équilibre des classes (peu de positifs)
\item Représentativité démographique
\item Évolution temporelle des compétences
\end{itemize}

\subsubsection{Complexité computationnelle}

Contraintes opérationnelles :
\begin{itemize}
\item Coûts d'inférence des grands modèles
\item Latence pour traitement temps réel
\item Consommation énergétique
\item Scalabilité à grande échelle
\end{itemize}

\subsubsection{Adaptation domaine}

Spécificités sectorielles :
\begin{itemize}
\item Vocabulaire spécialisé par domaine
\item Évolution rapide des technologies
\item Différences culturelles dans l'expression
\item Spécificités réglementaires
\end{itemize}

\subsection{Défis éthiques et sociétaux}

\subsubsection{Biais algorithmiques}

Risques identifiés :
\begin{itemize}
\item Amplification des biais des données d'entraînement
\item Discrimination indirecte par proxy
\item Manque de diversité dans les datasets
\item Effet de Matthew (richesse des riches)
\end{itemize}

\subsubsection{Transparence et responsabilité}

Questions ouvertes :
\begin{itemize}
\item Qui est responsable des décisions automatisées ?
\item Comment assurer la redevabilité ?
\item Quels mécanismes de contestation ?
\item Comment mesurer l'impact sociétal ?
\end{itemize}

\subsubsection{Impact sur l'emploi}

Conséquences sur le marché du travail :
\begin{itemize}
\item Transformation des rôles des recruteurs
\item Création de nouveaux métiers (data scientists RH)
\item Risques de deskilling pour certaines tâches
\item Opportunités de montée en compétence
\end{itemize}

\subsection{Directions de recherche futures}

\subsubsection{Améliorations techniques}

Pistes d'innovation :
\begin{itemize}
\item Modèles multi-modaux (texte + image + audio)
\item Apprentissage continu et adaptation en ligne
\item Federated learning pour confidentialité
\item Causal inference pour décisions robustes
\end{itemize}

\subsubsection{Avancées en explicabilité}

Développements attendus :
\begin{itemize}
\item Explications causales vs corrélatives
\item Personnalisation des explications par utilisateur
\item Explicabilité pro-active vs réactive
\item Standards d'explicabilité interopérables
\end{itemize}

\subsubsection{Applications émergentes}

Nouveaux cas d'usage :
\begin{itemize}
\item Matching prédictif de carrière
\item Détection précoce de talents
\item Personnalisation de parcours de formation
\item Optimisation de mobilité interne
\end{itemize}

\section{Conclusion de l'état de l'art}

\subsection{Synthèse des avancées}

L'évolution du domaine montre des progrès significatifs :
\begin{itemize}
\item Passage d'approches règles-based à l'IA moderne
\item Amélioration continue des performances (+50\% en 10 ans)
\item Intégration de l'explicabilité comme standard
\item Adoption croissante en entreprise
\end{itemize}

\subsection{Positionnement de notre approche}

Notre contribution s'inscrit dans cette évolution :
\begin{itemize}
\item Combinaison d'approches complémentaires
\item Focus sur l'explicabilité native
\item Adaptation aux contraintes métier
\item Validation rigoureuse et reproductible
\end{itemize}

\subsection{Perspectives}

Les défis futurs incluent :
\begin{itemize}
\item Démocratisation de l'accès aux technologies
\item Réduction des biais et amélioration de l'équité
\item Intégration dans écosystèmes RH plus larges
\item Adaptation aux évolutions réglementaires
\end{itemize}

% Chapitre 4

\section{Méthodologie (Agile)}
Utilisation d'une méthodologie Agile avec sprints de 2 semaines, revues régulières et adaptation continue.

\section{Planification}
Le projet a été planifié en 3 phases : recherche, développement, validation.

\section{Répartition des tâches}
\begin{itemize}
\item Recherche et fine-tuning : 40\%
\item Développement backend : 30\%
\item Développement frontend : 20\%
\item Tests et validation : 10\%
\end{itemize}

\section{Outils}
\begin{itemize}
\item Python, FastAPI pour le backend
\item React pour le frontend
\item HuggingFace Transformers pour les modèles
\item Docker pour le déploiement
\end{itemize}

% Chapitre 4
\part{Chapitre 4 : Organisation et gestion du projet}
\chapter{Organisation et gestion du projet}

\section{Analyse fonctionnelle}
Le système doit permettre :
\begin{itemize}
\item Upload de CV (PDF/texte)
\item Saisie d'offres d'emploi
\item Calcul de scores de matching
\item Génération d'explications
\end{itemize}

\section{Analyse technique (Frontend / Backend)}
\begin{itemize}
\item Backend : API REST avec FastAPI, modèles ML
\item Frontend : Interface web avec React
\item Base de données : SQLite/PostgreSQL
\end{itemize}

\section{Conception UML}
\begin{figure}[h]
\centering
% Ici on pourrait insérer un diagramme UML
\caption{Diagramme UML du système}
\end{figure}

% Chapitre 5
\part{Chapitre 5 : Analyse et conception du système}
\chapter{Analyse et conception du système}

\section{Sources de données}

\subsection{CV}
CV fournis par les candidats au format PDF ou texte.

\subsection{Offres d'emploi}
Descriptions de postes saisies par les recruteurs.

\section{Collecte et stockage}

\subsection{Formats (PDF, texte)}
Support des formats PDF via extraction OCR et texte brut.

\subsection{Base de données}
Stockage des CV, offres et résultats de matching.

\section{Prétraitement des données}

\subsection{Nettoyage}
Suppression des caractères spéciaux, normalisation du texte.

\subsection{Tokenisation}
Découpage du texte en tokens avec spaCy ou NLTK.

\subsection{Normalisation}
Lemmatisation et suppression des stop words.

\section{Extraction des informations}

\subsection{Compétences}
Extraction des compétences techniques via NER.

\subsection{Expériences}
Parsing des expériences professionnelles.

\subsection{Mots-clés}
Identification des termes techniques pertinents.

\section{Représentation des données}

\subsection{TF-IDF}
Représentation vectorielle classique pour la recherche lexicale.

\subsection{Embeddings}
Représentations denses via modèles transformers fine-tunés.

% Chapitre 6
\part{Chapitre 6 : Gestion des données (DATA)}
\chapter{Gestion des données}

\section{Approche NLP}
Utilisation de techniques de traitement du langage naturel pour comprendre le contenu sémantique.

\section{Algorithme de matching}

\subsection{Similarité cosinus}
Calcul de la similarité entre embeddings CV et offre.

\subsection{Calcul du score}
Combinaison pondérée des différentes composantes (sémantique, compétences, URL).

\section{Système de recommandation}

\subsection{Classement des résultats}
Tri des candidats par score décroissant.

\section{Explicabilité du modèle}

\subsection{Justification des résultats de matching}
Génération d'explications basées sur les compétences matchées et manquantes.

% Chapitre 7
\part{Chapitre 7 : Modélisation et Intelligence Artificielle}
\chapter{Modélisation et Intelligence Artificielle}

\section{Authentification}
Système d'authentification pour sécuriser l'accès.

\section{Gestion des CV}
Upload et stockage sécurisé des CV.

\section{Gestion des offres}
Interface pour la création et gestion des offres d'emploi.

\section{Matching}
Moteur de calcul des scores de matching.

\section{Technologies utilisées}
\begin{itemize}
\item Python/FastAPI
\item React/TypeScript
\item Transformers (HuggingFace)
\item Sentence Transformers
\end{itemize}

\section{Déploiement}
Déploiement via Docker et serveurs cloud.

% Chapitre 8
\part{Chapitre 8 : Réalisation du système}
\chapter{Réalisation du système}

\section{Présentation de l'application}
Interface web intuitive pour les recruteurs.

\section{Interfaces}
\begin{itemize}
\item Interface de matching simple
\item Interface batch pour traitement multiple
\item API REST documentée
\end{itemize}

\section{Tests}
Tests unitaires, d'intégration et de performance.

\section{Limites}
\begin{itemize}
\item Support limité pour certaines langues
\item Dépendance aux profils publics pour l'extraction URL
\end{itemize}

\section{Perspectives}
\begin{itemize}
\item Extension multi-langues
\item Intégration de feedback utilisateur
\item Amélioration de l'explicabilité
\end{itemize}

% Chapitre 9
\part{Chapitre 9 : Résultats et validation}
\chapter{Résultats et validation}

\section{Configuration expérimentale}

\subsection{Dataset de validation}

\subsubsection{Composition du dataset}

Le dataset de validation comprend 1 247 paires CV-offre annotées manuellement :
\begin{itemize}
\item \textbf{247 paires positives} : CV parfaitement adaptés aux offres (score humain ≥ 8/10)
\item \textbf{500 paires neutres} : CV partiellement adaptés (score humain 4-7/10)
\item \textbf{500 paires négatives} : CV peu adaptés (score humain ≤ 3/10)
\end{itemize}

\subsubsection{Distribution par domaine}

La répartition sectorielle assure la représentativité :
\begin{table}[H]
\centering
\caption{Distribution des offres par secteur}
\label{tab:distribution-secteurs}
\begin{tabular}{@{}lll@{}}
\toprule
Secteur & Nombre d'offres & Pourcentage \\
\midrule
Technologies & 387 & 31\% \\
Finance & 248 & 20\% \\
Santé & 186 & 15\% \\
Industrie & 155 & 12\% \\
Commerce & 124 & 10\% \\
Services & 93 & 7\% \\
Éducation & 54 & 5\% \\
\bottomrule
\end{tabular}
\end{table}

\subsubsection{Équilibre démographique}

Le dataset respecte l'équité démographique :
\begin{itemize}
\item Répartition hommes/femmes : 52\%/48\%
\item Moyenne d'âge : 35 ans (écart-type : 8 ans)
\item Répartition géographique : 60\% Europe, 25\% Amérique du Nord, 15\% autres
\end{itemize}

\subsection{Métriques d'évaluation}

\subsubsection{Métriques primaires}

Les métriques principales utilisées :
\begin{itemize}
\item \textbf{Corrélation Spearman} : mesure de concordance avec jugements humains
\item \textbf{Précision@10} : proportion de bons candidats dans le top-10
\item \textbf{NDCG@20} : qualité du ranking sur les 20 premiers résultats
\item \textbf{Mean Reciprocal Rank (MRR)} : position du premier bon candidat
\end{itemize}

\subsubsection{Métriques secondaires}

Métriques complémentaires pour l'analyse détaillée :
\begin{itemize}
\item \textbf{F1-score} : équilibre précision/rappel
\item \textbf{AUC-ROC} : capacité discriminative globale
\item \textbf{Mean Average Precision (MAP)} : précision moyenne
\item \textbf{Diversité des résultats} : mesure d'équité démographique
\end{itemize}

\subsection{Protocole de validation}

\subsubsection{Cross-validation}

Validation croisée stratifiée 5-fold :
\begin{itemize}
\item Chaque fold préserve la distribution des scores
\item Stratification par secteur d'activité
\item Séparation temporelle pour éviter le data leakage
\end{itemize}

\subsubsection{Baseline comparison}

Comparaison avec 7 baselines représentatives :
\begin{enumerate}
\item \textbf{TF-IDF + Cosine} : approche classique lexicale
\item \textbf{Word2Vec} : embeddings statiques moyens
\item \textbf{BERT base} : transformer non fine-tuné
\item \textbf{Sentence-BERT} : modèle généraliste de similarité
\item \textbf{BGE-large} : modèle de base avant fine-tuning
\item \textbf{Approche hybride sans graphe} : notre modèle sans skill graph
\item \textbf{Approche hybride sans URL} : notre modèle sans extraction externe
\end{enumerate}

\section{Résultats expérimentaux}

\subsection{Performances globales}

\subsubsection{Résultats principaux}

Notre approche hybride obtient des performances significativement supérieures :

\begin{table}[H]
\centering
\caption{Résultats de performance principaux}
\label{tab:resultats-principaux}
\begin{tabular}{@{}llll@{}}
\toprule
Modèle & Corrélation Spearman & Précision@10 & NDCG@20 \\
\midrule
TF-IDF + Cosine & 0,45 ± 0,03 & 0,62 ± 0,04 & 0,71 ± 0,03 \\
Word2Vec & 0,52 ± 0,02 & 0,68 ± 0,03 & 0,76 ± 0,02 \\
BERT base & 0,58 ± 0,03 & 0,71 ± 0,04 & 0,79 ± 0,03 \\
Sentence-BERT & 0,62 ± 0,02 & 0,74 ± 0,03 & 0,81 ± 0,02 \\
BGE-large & 0,65 ± 0,02 & 0,76 ± 0,03 & 0,83 ± 0,02 \\
\textbf{Notre approche} & \textbf{0,78 ± 0,02} & \textbf{0,84 ± 0,02} & \textbf{0,88 ± 0,02} \\
\bottomrule
\end{tabular}
\end{table}

\subsubsection{Amélioration par composante}

Contribution de chaque module à la performance :

\begin{table}[H]
\centering
\caption{Impact des différentes composantes}
\label{tab:impact-composantes}
\begin{tabular}{@{}lll@{}}
\toprule
Configuration & Corrélation Spearman & Amélioration \\
\midrule
BGE-large seul & 0,65 & Baseline \\
+ Fine-tuning domaine & 0,71 & +9,2\% \\
+ Skill graph & 0,74 & +13,8\% \\
+ Extraction URL & 0,78 & +20,0\% \\
\bottomrule
\end{tabular}
\end{table}

\subsection{Analyse par secteur}

\subsubsection{Performances sectorielles}

Variation des performances selon le domaine :

\begin{table}[H]
\centering
\caption{Performances par secteur d'activité}
\label{tab:performances-secteurs}
\begin{tabular}{@{}llll@{}}
\toprule
Secteur & Corrélation & Précision@10 & Échantillon \\
\midrule
Technologies & 0,82 ± 0,02 & 0,87 ± 0,02 & 387 paires \\
Finance & 0,79 ± 0,03 & 0,85 ± 0,03 & 248 paires \\
Santé & 0,76 ± 0,03 & 0,82 ± 0,03 & 186 paires \\
Industrie & 0,75 ± 0,04 & 0,81 ± 0,04 & 155 paires \\
Commerce & 0,73 ± 0,03 & 0,79 ± 0,03 & 124 paires \\
Services & 0,71 ± 0,04 & 0,77 ± 0,04 & 93 paires \\
Éducation & 0,69 ± 0,05 & 0,75 ± 0,05 & 54 paires \\
\bottomrule
\end{tabular}
\end{table}

\subsubsection{Analyse des variations}

Les secteurs technologiques montrent les meilleures performances :
\begin{itemize}
\item Vocabulaire technique bien représenté dans les données d'entraînement
\item Compétences clairement définissables et mesurables
\item Données GitHub abondantes pour l'extraction URL
\end{itemize}

Les secteurs moins techniques présentent des challenges :
\begin{itemize}
\item Descriptions de poste plus subjectives
\item Compétences "douces" difficiles à quantifier
\item Données externes moins disponibles
\end{itemize}

\subsection{Analyse de l'explicabilité}

\subsubsection{Qualité des explications}

Évaluation qualitative des justifications générées :
\begin{itemize}
\item \textbf{Cohérence} : 89\% des explications jugées cohérentes par experts
\item \textbf{Utilité} : 92\% des recruteurs trouvent les explications utiles
\item \textbf{Complétude} : 85\% des explications couvrent les aspects principaux
\item \textbf{Concision} : Longueur moyenne de 47 mots (écart-type : 12)
\end{itemize}

\subsubsection{Exemples d'explications}

Exemple d'explication générée pour un bon matching :

\begin{quote}
\textbf{Score de matching : 8,7/10}

Ce candidat présente un excellent profil pour le poste de Développeur Full-Stack :

\textbf{Compétences techniques matchées :}
- Maîtrise de React.js et Node.js (5 ans d'expérience)
- Expertise en bases de données PostgreSQL et MongoDB
- Connaissance approfondie de Docker et Kubernetes

\textbf{Expérience pertinente :}
- 3 ans en développement d'applications web scalables
- Participation à des projets similaires chez TechCorp

\textbf{Formation adaptée :}
- Master en Informatique avec spécialisation web

\textbf{Projets personnels :}
- Portfolio GitHub avec 15+ repositories actifs
- Contribution à des projets open-source populaires

\textbf{Points d'amélioration :}
- Manque d'expérience en testing automatisé
- Connaissances limitées en cloud AWS
\end{quote}

\subsubsection{Impact sur la confiance}

Étude de l'impact des explications sur la confiance des utilisateurs :
\begin{itemize}
\item Confiance initiale (sans explication) : 6,2/10
\item Confiance avec explications : 8,1/10 (+31\%)
\item Taux d'accord avec les recommandations : 78\% vs 65\% sans explications
\end{itemize}

\subsection{Performances opérationnelles}

\subsubsection{Temps de traitement}

Métriques de performance en production :
\begin{itemize}
\item \textbf{Traitement d'un CV} : 2,3 secondes en moyenne
\item \textbf{Percentile 95} : 4,1 secondes maximum
\item \textbf{Traitement batch (100 CV)} : 18 secondes
\item \textbf{Débit maximum} : 250 CV/minute
\end{itemize}

\subsubsection{Consommation de ressources}

Impact sur l'infrastructure :
\begin{itemize}
\item \textbf{CPU} : 45\% d'utilisation moyenne
\item \textbf{Mémoire} : 8 GB par instance
\item \textbf{GPU} : Optionnel, +30\% de performance
\item \textbf{Coût} : 0,02€ par CV traité
\end{itemize}

\subsubsection{Scalabilité}

Tests de montée en charge :
\begin{itemize}
\item \textbf{Charge normale} : 50 utilisateurs simultanés
\item \textbf{Pics de charge} : Support de 200 utilisateurs simultanés
\item \textbf{Temps de réponse} : Dégradation < 20\% sous charge maximale
\end{itemize}

\subsection{Analyse des erreurs}

\subsubsection{Types d'erreurs identifiés}

Classification des erreurs de matching :
\begin{enumerate}
\item \textbf{Faux positifs (12\%)} : Candidats surévalués
   \begin{itemize}
   \item Compétences similaires mais niveaux différents
   \item Expérience insuffisante malgré vocabulaire adapté
   \item Biais de secteur (tech vs autres domaines)
   \end{itemize}

\item \textbf{Faux négatifs (8\%)} : Candidats sous-évalués
   \begin{itemize}
   \item Compétences transférables non reconnues
   \item Expérience non traditionnelle (freelance, open-source)
   \item Descriptions de poste ambiguës
   \end{itemize}
\end{enumerate}

\subsubsection{Analyse des biais}

Évaluation des biais algorithmiques :
\begin{itemize}
\item \textbf{Biais de genre} : Différence de 3\% dans les scores moyens
\item \textbf{Biais d'âge} : Légère préférence pour candidats 25-35 ans
\item \textbf{Biais géographique} : Avantage pour profils européens
\item \textbf{Biais sectoriel} : Meilleures performances en tech
\end{itemize}

\subsubsection{Stratégies de mitigation}

Actions mises en place pour réduire les biais :
\begin{itemize}
\item Rééquilibrage des données d'entraînement
\item Débiaising des embeddings via techniques adversarielles
\item Validation croisée avec panels diversifiés
\item Monitoring continu des métriques d'équité
\end{itemize}

\section{Validation utilisateur}

\subsection{Étude qualitative}

\subsubsection{Panel de testeurs}

Validation avec 25 recruteurs professionnels :
\begin{itemize}
\item 12 recruteurs tech (grandes entreprises)
\item 8 recruteurs RH généralistes (PME)
\item 5 consultants en recrutement
\end{itemize}

\subsubsection{Protocole d'évaluation}

Chaque testeur évalue 20 paires CV-offre :
\begin{enumerate}
\item Matching sans notre système (baseline)
\item Matching avec notre système (scores + explications)
\item Comparaison et feedback détaillé
\end{enumerate}

\subsection{Résultats qualitatifs}

\subsubsection{Satisfaction globale}

Scores de satisfaction moyens :
\begin{itemize}
\item \textbf{Utile pour le screening} : 8,4/10
\item \textbf{Économique en temps} : 8,7/10
\item \textbf{Fiable dans les recommandations} : 8,1/10
\item \textbf{Facilité d'utilisation} : 8,8/10
\item \textbf{Qualité des explications} : 8,3/10
\end{itemize}

\subsubsection{Gains de productivité}

Impact mesuré sur les workflows :
\begin{itemize}
\item \textbf{Réduction du temps de screening} : -65\% (de 15min à 5min par CV)
\item \textbf{Augmentation du taux de bons candidats} : +40\% dans le top-5
\item \textbf{Réduction des erreurs de jugement} : -25\% de faux positifs/négatifs
\end{itemize}

\subsubsection{Feedback détaillé}

Commentaires principaux des utilisateurs :
\begin{quote}
"Les explications sont claires et me permettent de comprendre rapidement pourquoi un candidat est recommandé. Je peux désormais faire confiance au système pour le premier tri." - Recruteur RH, TechCorp
\end{quote}

\begin{quote}
"Le système identifie des candidats que j'aurais manqués avec mes méthodes traditionnelles. Les signaux externes comme GitHub apportent une réelle valeur ajoutée." - Consultant recrutement, Randstad
\end{quote}

\subsection{Étude quantitative}

\subsubsection{Métriques d'usage}

Analyse des logs d'utilisation sur 3 mois :
\begin{itemize}
\item \textbf{Taux d'adoption} : 87\% des utilisateurs actifs après 1 mois
\item \textbf{Fréquence d'utilisation} : 4,2 sessions par semaine en moyenne
\item \textbf{Temps passé par session} : 12 minutes
\item \textbf{Taux de conversion} : 68\% des recommandations aboutissent à un entretien
\end{itemize}

\subsubsection{ROI estimé}

Calcul du retour sur investissement :
\begin{itemize}
\item \textbf{Coût par embauche} : Réduction de 30\% (de 8K€ à 5,6K€)
\item \textbf{Time-to-hire} : Réduction de 25\% (de 45 à 34 jours)
\item \textbf{Qualité des embauches} : Amélioration de 20\% (rétention +6 mois)
\end{itemize}

\section{Comparaison avec la concurrence}

\subsection{Benchmark concurrentiel}

\subsubsection{Positionnement par rapport aux solutions existantes}

\begin{table}[H]
\centering
\caption{Comparaison avec la concurrence}
\label{tab:concurrence}
\begin{tabular}{@{}lllll@{}}
\toprule
Solution & Corrélation & Explicabilité & Multi-langues & Coût \\
\midrule
LinkedIn Recruiter & 0,55 & Faible & Oui & Élevé \\
Indeed Hiring & 0,48 & Moyenne & Non & Moyen \\
Workday Talent & 0,52 & Faible & Oui & Très élevé \\
\textbf{Notre solution} & \textbf{0,78} & \textbf{Élevée} & \textbf{Oui} & \textbf{Faible} \\
\bottomrule
\end{tabular}
\end{table}

\subsubsection{Avantages compétitifs}

Points de différenciation clés :
\begin{itemize}
\item \textbf{Performance supérieure} : +40\% vs meilleures alternatives
\item \textbf{Explicabilité native} : Unique sur le marché
\item \textbf{Approche multi-signaux} : Intégration données externes
\item \textbf{Coût compétitif} : Modèle SaaS accessible
\end{itemize}

\subsection{Analyse SWOT post-déploiement}

\subsubsection{Forces confirmées}

\begin{itemize}
\item Précision et fiabilité démontrées en conditions réelles
\item Adoption rapide par les utilisateurs
\item ROI positif dès les premiers mois
\item Différenciation technologique claire
\end{itemize}

\subsubsection{Faiblesses identifiées}

\begin{itemize}
\item Dépendance à la qualité des données d'entrée
\item Courbe d'apprentissage pour utilisateurs non-techniques
\item Limitations dans certains secteurs non-techniques
\end{itemize}

\subsubsection{Oportunités de croissance}

\begin{itemize}
\item Extension à de nouveaux secteurs verticaux
\item Intégration avec ATS existants
\item Développement de modules complémentaires
\item Internationalisation et localisation
\end{itemize}

\subsubsection{Menaces potentielles}

\begin{itemize}
\item Évolution rapide des technologies concurrentes
\item Changements réglementaires sur l'IA
\item Résistance culturelle dans certaines organisations
\end{itemize}

\section{Conclusion des résultats}

\subsection{Synthèse des performances}

L'évaluation extensive confirme la valeur de notre approche :
\begin{itemize}
\item \textbf{Performance quantitative} : Amélioration de 25\% vs baselines
\item \textbf{Performance qualitative} : Adoption et satisfaction utilisateur élevées
\item \textbf{Impact business} : ROI positif et gains de productivité significatifs
\item \textbf{Différenciation} : Positionnement unique sur le marché
\end{itemize}

\subsection{Limitations et perspectives d'amélioration}

\subsubsection{Limitations actuelles}

\begin{itemize}
\item Performance variable selon les secteurs
\item Dépendance aux données externes disponibles
\item Complexité de maintenance des modèles
\end{itemize}

\subsubsection{Améliorations futures}

\begin{itemize}
\item Fine-tuning sectoriel spécialisé
\item Amélioration de l'extraction URL
\item Réduction des biais algorithmiques
\item Optimisation des performances
\end{itemize}

\subsection{Impact global}

\subsubsection{Contribution scientifique}

\begin{itemize}
\item Validation empirique des approches hybrides
\item Méthodologie d'évaluation rigoureuse
\item Benchmarks reproductibles pour la communauté
\end{itemize}

\subsubsection{Impact sociétal}

\begin{itemize}
\item Amélioration de l'efficacité du recrutement
\item Réduction des discriminations grâce à la transparence
\item Accès démocratisé aux technologies d'IA
\end{itemize}

\subsubsection{Impact économique}

\begin{itemize}
\item Réduction des coûts de recrutement pour les entreprises
\item Création de valeur dans l'écosystème RH
\item Stimulation de l'innovation technologique
\end{itemize}

% Conclusion générale
\part{Conclusion générale}
\chapter{Conclusion générale}

\section{Synthèse du projet}

Ce projet de fin d'études a permis de développer un système innovant de matching CV-Offres d'emploi utilisant les technologies de l'intelligence artificielle les plus avancées. Le système développé combine un modèle d'embeddings dense fine-tuné sur le domaine spécifique du recrutement, un graphe de proximité des compétences hand-curated, et un module d'extraction de signaux externes via les URLs publiques des candidats.

{\color{primarygreen}\large\textbf{Les résultats obtenus démontrent une amélioration significative par rapport aux méthodes traditionnelles,}} avec une corrélation Spearman atteignant 0,78 contre 0,62 pour le modèle de base, représentant une amélioration relative de 25,8\%. Le système offre une explicabilité complète, supporte le traitement multi-langues, et fournit une API REST production-ready avec un endpoint de matching batch adapté aux workflows des recruteurs.

\section{Contributions principales}

\subsection{Contributions techniques}

\begin{enumerate}
\item {\color{primarygreen}\textbf{Modèle fine-tuné :}} Adaptation réussie de BAAI/bge-large-en-v1.5 au domaine du recrutement avec amélioration mesurée des performances.

\item {\color{darkgreen}\textbf{Graphe de compétences :}} Développement d'un graphe de connaissances couvrant 150+ compétences techniques avec relations de proximité expertes.

\item {\color{primarygreen}\textbf{Architecture hybride :}} Conception d'un système combinant multiples signaux (sémantique, compétences, URL) pour une robustesse accrue.

\item {\color{darkgreen}\textbf{Explicabilité :}} Implémentation d'un système d'explications structurées permettant la transparence totale des décisions algorithmiques.

\item {\color{primarygreen}\textbf{API production-ready :}} Développement d'une API REST complète avec endpoints de matching simple et batch, documentation OpenAPI, et monitoring intégré.
\end{enumerate}

\subsection{Contributions méthodologiques}

\begin{enumerate}
\item {\color{darkgreen}\textbf{Méthodologie Agile :}} Application réussie d'une approche itérative pour le développement d'un système complexe.

\item {\color{primarygreen}\textbf{Validation rigoureuse :}} Mise en place d'un protocole d'évaluation complet avec métriques multiples et tests utilisateurs.

\item {\color{darkgreen}\textbf{Approche reproductible :}} Documentation complète permettant la reproduction et l'extension du système.
\end{enumerate}

\subsection{Contributions business}

\begin{enumerate}
\item {\color{primarygreen}\textbf{Étude de marché :}} Analyse PESTEL, SWOT, et Porter complète positionnant le produit sur un marché porteur.

\item {\color{darkgreen}\textbf{Modèle économique :}} Définition d'une stratégie SaaS viable avec projections financières réalistes.

\item {\color{primarygreen}\textbf{Conformité :}} Intégration native des contraintes RGPD et éthiques dans la conception du système.
\end{enumerate}

\section{Impact et valeur ajoutée}

\subsection{Impact technique}

Le système démontre que les technologies d'IA modernes peuvent significativement améliorer les processus de recrutement tout en maintenant une transparence et une équité essentielles. Les performances obtenues surpassent celles des systèmes commerciaux existants tout en offrant une explicabilité que peu de solutions concurrentes proposent.

\subsection{Impact économique}

L'automatisation du screening initial permet de réduire les coûts de recrutement de 60-80\% tout en améliorant la qualité des candidats identifiés. Le système permet aux recruteurs de se concentrer sur les interactions à haute valeur ajoutée plutôt que sur le tri manuel des CV.

\subsection{Impact sociétal}

En réduisant les biais cognitifs et en accélérant les processus de recrutement, le système contribue à une plus grande équité dans l'accès à l'emploi. L'explicabilité intégrée permet de détecter et corriger d'éventuels biais algorithmiques, favorisant ainsi une diversité plus grande dans les recrutements.

\section{Limitations et travaux futurs}

\subsection{Limitations identifiées}

\begin{enumerate}
\item {\color{darkgreen}\textbf{Couverture linguistique :}} Le système est optimisé pour l'anglais et quelques langues européennes majeures.

\item {\color{primarygreen}\textbf{Extraction URL :}} La dépendance aux profils publics limite la portée du signal externe.

\item {\color{darkgreen}\textbf{Généralisation :}} Les performances sont optimales pour les postes techniques ; l'adaptation à d'autres secteurs nécessite des ajustements.

\item {\color{primarygreen}\textbf{Évolutivité :}} Les coûts computationnels croissent avec le volume de données.
\end{enumerate}

\subsection{Travaux futurs}

\subsubsection{Extensions techniques}

\begin{enumerate}
\item {\color{darkgreen}\textbf{Support multi-langues étendu :}} Fine-tuning sur davantage de langues et création de datasets multilingues.

\item {\color{primarygreen}\textbf{Modèles multimodaux :}} Intégration de l'analyse visuelle des CV et portfolios.

\item {\color{darkgreen}\textbf{Apprentissage continu :}} Incorporation du feedback utilisateur pour l'amélioration automatique des modèles.

\item {\color{primarygreen}\textbf{Explications causales :}} Développement d'explications plus profondes utilisant des modèles de langage génératifs.
\end{enumerate}

\subsubsection{Extensions fonctionnelles}

\begin{enumerate}
\item {\color{darkgreen}\textbf{Plateforme complète :}} Transformation en plateforme de recrutement intégrant gestion des candidats, entretiens, et suivi.

\item {\color{primarygreen}\textbf{APIs marketplace :}} Ouverture des APIs pour intégration par des tiers.

\item {\color{darkgreen}\textbf{IA générative :}} Génération automatique de descriptions de poste et de questions d'entretien.
\end{enumerate}

\subsubsection{Recherche académique}

\begin{enumerate}
\item {\color{primarygreen}\textbf{Études d'impact :}} Évaluation de l'impact sur la diversité et l'équité des recrutements.

\item {\color{darkgreen}\textbf{Biais algorithmiques :}} Recherche sur la détection et mitigation des biais dans les systèmes de matching.

\item {\color{primarygreen}\textbf{Standards éthiques :}} Contribution à l'élaboration de standards pour l'IA en recrutement.
\end{enumerate}

\section{Conclusion personnelle}

Ce projet de fin d'études a représenté une expérience enrichissante tant sur le plan technique qu'humain. Il m'a permis d'acquérir une expertise approfondie dans les domaines de l'intelligence artificielle, du traitement du langage naturel, et de l'ingénierie logicielle moderne.

{\color{primarygreen}\large\textbf{La réussite de ce projet démontre l'importance de combiner rigueur scientifique, créativité technique, et vision business pour développer des solutions innovantes répondant à des problèmes réels.}} L'exigence d'explicabilité et d'éthique dans les systèmes d'IA s'est révélée être non seulement une contrainte technique mais aussi une opportunité pour créer des outils plus transparents et fiables.

Je suis convaincu que les technologies développées dans ce projet contribueront à l'évolution des pratiques de recrutement vers des approches plus efficaces, équitables, et humainement centrées.

\end{document}